\section{Related Work}
\label{sec:related}

\subsection{Stippling and Point-Based Non-Photorealistic Rendering}

Computational stippling traces its origins to the floating-points
method of \citet{deussen2000floating}, which used Lloyd relaxation
\citep{lloyd1982least} to distribute points according to a target
density field. \citet{secord2002weighted} advanced this with Weighted
Voronoi Stippling, combining rejection sampling, density-weighted
centroid computation, and Lloyd relaxation to achieve blue-noise
point distributions; this remains the gold standard and forms our
algorithmic foundation. Subsequent work explored capacity-constrained
distributions \citep{balzer2009capacity}, example-based texture
synthesis \citep{kim2008stippling}, and perceptual studies of dot
quality \citep{spicker2019quantifying}. Poisson disk sampling
\citep{bridson2007fast} offers an $O(n)$ alternative but does not
directly support density modulation.

A key limitation of all existing stippling methods is that dot size
and density are mapped to a single scalar field---typically luminance.
No prior work, to our knowledge, has incorporated scene depth as an
additional dimension for dot size modulation. Our depth-to-radius
mapping addresses this gap by treating depth as a first-class
artistic control variable.

\subsection{Monocular Depth Estimation}

Monocular depth estimation---predicting per-pixel depth from a single
RGB image---has undergone rapid advancement through deep learning.
Early CNN-based methods \citep{eigen2014depth} demonstrated the
feasibility of single-image depth regression. Subsequent architectures
incorporated skip connections \citep{laina2016deeper}, ordinal
regression \citep{fu2018deep}, and transformer backbones
\citep{ranftl2021vision} to improve accuracy and generalization.

The YOLO26 model family \citep{ultralytics2025yolo26} represents the
latest generation of unified real-time vision models, supporting
multiple tasks---including depth estimation, semantic segmentation,
object detection, and pose estimation---within a single architecture.
The depth estimation variant produces dense metric depth maps (in
meters) at real-time speeds on consumer GPUs, while the semantic
segmentation variant assigns per-pixel class labels from the Cityscapes
19-class taxonomy \citep{cordts2016cityscapes}. The availability of
both depth and semantic information from a single model family
simplifies pipeline integration and ensures geometric and semantic
consistency.

In the context of non-photorealistic rendering, depth maps have been
used primarily for view-dependent effects such as depth-of-field
simulation and atmospheric perspective. \citet{winnemoller2006real}
employed depth for stylized focus effects. However, the use of depth
to modulate the size of individual rendering primitives (dots, strokes)
in a stippling or pointillism context is novel.

\subsection{Neural Style Transfer and Diffusion-Based Stylization}

Neural style transfer (NST), introduced by \citet{gatys2015neural},
separates content and style features using deep network activations
and recombines them via Gram matrix optimization. \citet{johnson2016perceptual}
accelerated this process with feed-forward transformer networks, and
\citet{ulyanov2016instance} identified instance normalization as a
critical component for fast stylization. While NST can produce
painterly effects, it does not provide explicit control over
individual rendering primitives (dots, strokes) and tends to produce
smooth, texture-mapped results rather than discrete dot patterns.

Diffusion-based methods offer higher-quality synthesis at the cost of
greater computational demands. \citet{rombach2022high} introduced
Latent Diffusion Models (Stable Diffusion), and \citet{zhang2023adding}
proposed ControlNet for conditional control via structural inputs
(Canny edges, depth maps, poses). \citet{hu2022lora} developed
low-rank adaptation (LoRA) for efficient fine-tuning, enabling
style-specific model customization with limited training data. A
Kusama-style LoRA could be trained on her body of work and combined
with a depth-conditioned ControlNet to generate polka dot art.

Our pipeline occupies a complementary niche: rather than relying on
learned style transfer, we use explicit geometric algorithms (Voronoi
stippling) with depth-driven parameter modulation. This provides
several advantages: (i) full interpretability---every dot's position,
size, and color is determined by a known mathematical rule; (ii)
controllability---artists can adjust parameters ($\alpha$, $\sigma$,
palette) to achieve specific effects; and (iii) determinism with
controlled randomness---the random seed is reproducible while the
Gaussian noise term provides organic variation. The output of our
pipeline could serve as a conditioning image for diffusion-based
refinement, combining the structural precision of algorithmic
stippling with the textural richness of learned style transfer.

\subsection{Art-Style Rendering}

Computational rendering of specific artistic styles has been explored
for impressionism \citep{hertzmann1998painterly}, pointillism
\citep{litwinowicz1997processing}, watercolor
\citep{curtis1997computer}, and other media. These methods typically
employ physics-based or procedural models of the medium's
characteristic marks. Kusama's polka dot art, despite its iconic
status, has received relatively little attention in the computational
rendering literature. Our work contributes a domain-specific framework
that encodes Kusama's four signature visual traits---hard-edged dots,
multi-scale coexistence, high-chroma palettes, and horror vacui---as
algorithmic parameters.
